%! Parametrization of Implicitly Defined Functions — Unit 4, Lesson 4.7 — Exit Ticket (Answer Key)
\documentclass[10pt]{article}
\usepackage{precalculus-article}
\usepackage{precalculus-key}   % pulls in -boxes; do NOT also load -boxes
\newcommand{\TallMath}[1]{$\displaystyle #1\rule[-1.4em]{0pt}{3.2em}$}

\begin{document}

\pageheader{Unit 4, Lesson 4.7}{Exit Ticket --- Answer Key}
\namedateperiod

\vspace{0.12in}

\begin{notesbox}{Parametrization Check --- Key}

The ellipse $\dfrac{(x+1)^2}{9} + \dfrac{(y-2)^2}{16} = 1$ is implicitly defined.

\vspace{10pt}
\textbf{1.}\quad Identify $h$, $k$, $a$, and $b$. Write parametric equations.

\vspace{6pt}
$h = $ \ans{$-1$} \quad $k = $ \ans{$2$} \quad $a = $ \ans{$3$} \quad $b = $ \ans{$4$}

\vspace{10pt}
$x(t) = $ \ans{$-1 + 3\cos t$} \qquad $y(t) = $ \ans{$2 + 4\sin t$}

\begin{teachernote}
  A common error: students write $x(t)=h\cos t$ instead of $x(t)=h+a\cos t$ ---
  forgetting to add the center coordinates. Prompt: ``The parametrization shifts the
  unit ellipse to be centered at $(h,k)$. What is the center here?''
\end{teachernote}

\vspace{10pt}
\textbf{2.}\quad Topmost and bottommost points.

\vspace{6pt}
\ans{$y(t) = 2 + 4\sin t$ is maximized when $\sin t = 1$, i.e., $t = \pi/2$.\\[3pt]
At $t = \pi/2$: $x = -1 + 3\cos(\pi/2) = -1 + 0 = -1$, $y = 2 + 4(1) = 6$.
\textbf{Topmost point: $(-1,\,6)$.}\\[3pt]
$y(t)$ is minimized when $\sin t = -1$, i.e., $t = 3\pi/2$.\\[3pt]
At $t = 3\pi/2$: $x = -1 + 3(0) = -1$, $y = 2 + 4(-1) = -2$.
\textbf{Bottommost point: $(-1,\,-2)$.}}

\end{notesbox}

\end{document}
